% !TEX root = task-1-17-1.tex
% Editable task bank. Append new BankTask blocks at the end of the document.
\documentclass[12pt,a4paper]{article}
\usepackage{physics-textbook}
\usepackage{graphicx,needspace}
\usetikzlibrary{decorations.pathmorphing}
\newif\ifShowSolutions
\ifdefined\TutorEdition
  \ShowSolutionstrue
\else
  \ShowSolutionsfalse % Change to \ShowSolutionstrue to display solutions.
\fi
\newcounter{banktask}
\newsavebox{\BankTaskBox}
\newcommand{\BankTaskContent}[3]{%
  \noindent\textbf{\thebanktask.}\quad #2\par
  \ifShowSolutions
    \Needspace{50mm}% Keep the solution heading with its opening text or diagram.
    #3\par
    \ifdefined\TutorEdition
      \ifcsname TutorSource#1\endcsname
        \nopagebreak\smallskip{\small\csname TutorSource#1\endcsname\par}%
      \fi
    \fi
  \fi
}
\newcommand{\BankTask}[3]{%
  \refstepcounter{banktask}\par\medskip
  \sbox{\BankTaskBox}{\begin{minipage}{\linewidth}%
    \BankTaskContent{#1}{#2}{#3}\end{minipage}}%
  \ifdim\dimexpr\ht\BankTaskBox+\dp\BankTaskBox\relax>\dimexpr\textheight-20mm\relax
    \Needspace{35mm}\BankTaskContent{#1}{#2}{#3}%
  \else
    \Needspace{\dimexpr\ht\BankTaskBox+\dp\BankTaskBox\relax}%
    \noindent\usebox{\BankTaskBox}\par
  \fi
  \medskip
}
\begin{document}
\BookSubsection{1.17.1}{Гармонические колебания}


% Source sdamgia-631: https://phys-ege.sdamgia.ru/problem?id=631
\BankTask{sdamgia-631}{
Скорость тела, совершающего гармонические колебания, меняется с течением времени в соответствии с уравнением $v =3 \cdot 10 ^{ - 2 } \sin 2 \pi t,$ где все величины выражены в СИ. Какова амплитуда колебаний скорости?
}{
\smallskip\noindent\textbf{Решение}\par
Общий вид закона изменения скорости тела со временем, совершающего колебания, имеет вид:

$v = v _{\max} \sin ( \omega _0t + \varphi _0 ) ,$

где $v _{\max}$  — амплитуда колебаний скорости. 

Сравнивая с $v =3 \cdot 10 ^{ - 2 } \sin 2 \pi t,$ заключаем, что амплитуда колебаний скорости равна $v _{\max} =3 \cdot 10 ^{ - 2 } \,\text{м/с}=0{,}03\,\text{м/с}.$

Ответ: $0{,}03$.
}

% Source sdamgia-3788: https://phys-ege.sdamgia.ru/problem?id=3788
\BankTask{sdamgia-3788}{
В таблице представлены данные о положении шарика, гармонически колеблющегося вдоль оси $Ox$ в различные моменты времени.

\begin{center}\begin{tabular}{|c|c|c|c|c|c|c|c|c|c|}\hline
$t$, с&$0$&$0{,}2$&$0{,}4$&$0{,}6$&$0{,}8$&$1$&$1{,}2$&$1{,}4$&$1{,}6$\\\hline
$x$, мм&$0$&$2$&$5$&$10$&$13$&$15$&$13$&$10$&$5$\\\hline
\end{tabular}\end{center}
\begin{center}\begin{tabular}{|c|c|c|c|c|c|c|c|c|}\hline
$t$, с&$1{,}8$&$2$&$2{,}2$&$2{,}4$&$2{,}6$&$2{,}8$&$3$&$3{,}2$\\\hline
$x$, мм&$2$&$0$&$-2$&$-5$&$-10$&$-13$&$-15$&$-13$\\\hline
\end{tabular}\end{center}

Какова амплитуда колебаний шарика?
}{
\smallskip\noindent\textbf{Решение}\par
Колебания симметричны относительно положения равновесия. Амплитудой называется величина максимального отклонения из положения равновесия. Из приведенной таблицы видно, что колебания симметричны относительно точки $x=0\,\text{мм},$ а максимальное отклонение достигается впервые через $1{,}0\,\text{с}$ и составляет $15\,\text{мм}$, именно такова величина амплитуды данного колебания.

Ответ: $15$.
}

% Source sdamgia-609: https://phys-ege.sdamgia.ru/problem?id=609
\BankTask{sdamgia-609}{
На рисунке дан график зависимости координаты материальной точки от времени. Какова частота колебаний?
\begin{center}\begin{tikzpicture}[task graph,x=.7cm,y=.36cm]
\draw[task grid,xstep=.5,ystep=1] (0,-4) grid (7,4);\draw[task axis] (-.4,0)--(7.3,0) node[right] {$t$, с};\draw[task axis] (0,-4)--(0,4.3) node[above] {$x$, см};
\draw[task curve,domain=0:6.3,samples=100] plot (\x,{3*sin(90*\x)});
\foreach \x in {2,4,6}{\node[below] at (\x,0) {$\x$};}\foreach \y in {-2,2}{\node[left] at (0,\y) {$\y$};}\node[below left] at (0,0) {$0$};
\end{tikzpicture}\end{center}
}{
\smallskip\noindent\textbf{Решение}\par
Из графика видно, что период колебания тела составляет $4$ секунды. Отсюда находим частоту:

$\nu = \frac{1}{T} = \frac{1}{4\,\text{с}} =0{,}25\,\text{Гц}.$

Ответ: $0{,}25$.
}

% Source sdamgia-604: https://phys-ege.sdamgia.ru/problem?id=604
\BankTask{sdamgia-604}{
Груз колеблется на пружине, подвешенной вертикально к потолку, при этом максимальное расстояние от потолка до центра груза равно $H$, минимальное $h$. Положение равновесия груза находится от потолка на расстоянии:

\noindent 1) $h$ \par
\noindent 2) $H$ \par
\noindent 3) $\frac{H + h}{2}$ \par
\noindent 4) $\frac{H - h}{2}$\par
}{
\smallskip\noindent\textbf{Решение}\par
Пружинный маятник совершает гармонические колебания вокруг своего положения равновесия. Пусть $A$  — амплитуда колебания, а $x$  — искомое расстояние от потолка до положения равновесия. Тогда 

$A=x - h=H - x,$

 откуда находим, что 

$x= \frac{H + h}{2} .$

Ответ: $3$.
}

% Source sdamgia-719: https://phys-ege.sdamgia.ru/problem?id=719
\BankTask{sdamgia-719}{
При свободных колебаниях груза на нити математического маятника его кинетическая энергия изменяется от $0\,\text{Дж}$ до $50\,\text{Дж}$, максимальное значение потенциальной энергии $50\,\text{Дж}$. Чему равна полная механическая энергия груза при таких колебаниях?
}{
\smallskip\noindent\textbf{Решение}\par
При свободных колебаниях груза на нити выполняется закон сохранения полной механической энергии:

$E = E_{\,\text{кин}} ( t ) + E_{\,\text{пот}} ( t ) =const.$

Минимуму кинетической энергии соответствует максимум потенциальной, и наоборот. Таким образом, полная механическая энергия груза во время колебаний равна $E=0\,\text{Дж} + 50\,\text{Дж}=50\,\text{Дж}.$

Ответ: $50$.
}

% Source sdamgia-409: https://phys-ege.sdamgia.ru/problem?id=409
\BankTask{sdamgia-409}{
Груз массой $m$ на пружине, совершая свободные колебания, проходит положение равновесия со скоростью $v .$ Через четверть периода колебаний он достигает положения максимального удаления от положения равновесия. Чему равен модуль изменения кинетической энергии груза за это время?

\noindent 1) $m v ^2$ \par
\noindent 2) $2m v ^2$ \par
\noindent 3) $\frac{m v ^2}{2}$ \par
\noindent 4) $0$\par
}{
\smallskip\noindent\textbf{Решение}\par
Через четверть периода, когда маятник достигает положения максимального отклонения, его скорость обращается в ноль. Таким образом, модуль изменения кинетической энергии за это время равен 

$| 0 - \frac{m v ^2}{2} |= \frac{m v ^2}{2} .$

Ответ: $3$.
}

% Source sdamgia-716: https://phys-ege.sdamgia.ru/problem?id=716
\BankTask{sdamgia-716}{
Груз массой $m$ свободно колеблется на пружине в горизонтальной плоскости. В начальный момент времени он проходит положение равновесия со скоростью $v ,$ и через четверть периода колебаний достигает положения максимального удаления. Модуль изменения полной механической энергии груза за это время равен

\noindent 1) $0$ \par
\noindent 2) $\frac{mv^2}{2}$ \par
\noindent 3) $m v ^2$ \par
\noindent 4) $2m v ^2$\par
}{
\smallskip\noindent\textbf{Решение}\par
При решении данной задачи необходимо обратить внимание, что нужно найти изменение энергии для груза, а не для системы груз  — пружина. Полная механическая энергия сохраняется только в замкнутой системе или в системе, на которую действуют силы, не совершающие работу. На груз действует пружина, которая совершает работу по его разгону и торможению, поэтому полная энергия груза изменяется.
Поскольку груз колеблется в горизонтальной плоскости, его потенциальная энергия не изменяется. Таким образом, изменение полной механической энергии груза будет равно изменению его кинетической энергии. В крайнем положении груз имеет нулевую скорость. Следовательно, изменение энергии равно $\frac{m v ^2}{2} .$

Ответ: $2$.

Примечание. Обратите внимание, что если бы груз колебался в вертикальной плоскости, то менялась бы не только его кинетическая, но и потенциальная энергия $mgh.$ А значит, ответ в задаче был бы другой. При этом было бы важно, в каком направлении за четверть периода сместился груз: в нижнее или в верхнее. Изменение потенциальной энергии в этих двух случаях будет разным.
}

% Source sdamgia-731: https://phys-ege.sdamgia.ru/problem?id=731
\BankTask{sdamgia-731}{
Подвешенный на нити грузик совершает гармонические колебания. В таблице представлены координаты грузика через одинаковые промежутки времени. Какова примерно максимальная скорость грузика? 

\begin{center}\begin{tabular}{|c|c|c|c|c|c|c|c|c|}\hline
$t$, $\text{с}$&$0$&$0{,}5$&$1$&$1{,}5$&$2$&$2{,}5$&$3$&$3{,}5$\\\hline
$x$, $\text{см}$&$6$&$3$&$0$&$3$&$6$&$3$&$0$&$3$\\\hline
\end{tabular}\end{center}
}{
\smallskip\noindent\textbf{Решение}\par
Максимальная скорость груза маятника связана с амплитудой колебаний и циклической частотой соотношением: 

$v _{\,\text{м}}=x_{\,\text{м}} \omega_0.$

Из таблицы видно, что при колебаниях смещение меняется от $0$ до $6\,\text{см}$. Следовательно, в положении равновесия смещение равно $3\,\text{см}$, а амплитуда колебаний равна $x_{\,\text{м}}=3\,\text{см},$ период $T=2\,\text{с},$ а значит, циклическая частота:

$\omega_0= \frac{2 \pi}{T} = \frac{2 \pi}{2\,\text{с}} = \pi \,\text{рад/с}.$

Таким образом, максимальная скорость груза равна приблизительно:

$v_{\text{м}}=0{,}03\,\text{м}\cdot3{,}14\,\text{с}^{-1}\approx0{,}09\,\text{м/с}.$

Ответ: $0{,}09\,\text{м/с}$.
}

% Source sdamgia-29740: https://phys-ege.sdamgia.ru/problem?id=29740
\BankTask{sdamgia-29740}{
В таблице представлена зависимость координаты тела от времени для пружинного маятника.

\begin{center}\begin{tabular}{|c|c|c|c|c|c|c|c|c|c|}\hline
$t$, $\text{с}$&$0$&$0{,}2$&$0{,}4$&$0{,}6$&$0{,}8$&$1$&$1{,}2$&$1{,}4$&$1{,}6$\\\hline
$x$, $\text{см}$&$20$&$14{,}2$&$0$&$-14{,}2$&$-20$&$-14{,}2$&$0$&$14{,}2$&$20$\\\hline
\end{tabular}\end{center}

Определите максимальную скорость маятника в процессе движения.
}{
\smallskip\noindent\textbf{Решение}\par
Из таблицы, представленной в условии, находим, что амплитуда колебаний равна $A=20 \,\text{см}=0{,}2 \,\text{м},$ а период колебаний $T=1{,}6 \,\text{с}.$ Максимальная скорость при колебательном движении $v _{\max} =\omega A,$ где циклическая частота $\omega= \frac{2 \pi}{T} .$ Тогда максимальная скорость движения: 

$v _{\max} = \frac{2 \pi A}{T} = \frac{2 \cdot 3{,}14 \cdot 0{,}2}{1{,}6} =0{,}785 \,\text{м/с}.$ 

Ответ: $0{,}785\,\text{м}/\text{с}$.
}

% Source sdamgia-50844: https://phys-ege.sdamgia.ru/problem?id=50844
\BankTask{sdamgia-50844}{
Груз массой $2\,\text{кг}$, закреплённый на пружине жесткостью $200\,\text{Н/м}$, совершает горизонтальные гармонические колебания. Максимальное ускорение груза при этом равно $10\,\text{м/с}^{2}$. Какова максимальная скорость груза?
}{
\smallskip\noindent\textbf{Решение}\par
При гармонических колебаниях максимальная сила $F_{\max}=kA,$ где $A$  — амплитуда колебаний. По второму закону Ньютона $F_{\max}=ma_{\max}.$ Тогда амплитуда колебаний равна

$A = \frac{ma_{\max}}{k} .$

Для гармонических колебаний выполняется закон сохранения энергии:

$\frac{kA ^2}{2} = \frac{m v _{\max} ^2}{2} .$

Объединяя уравнения, находим максимальную скорость груза:

$v _{\max}=a_{\max} \sqrt{\frac{m}{k}} =10 \cdot \sqrt{\frac{2}{200}} =1 \,\text{м/с}.$

Ответ: $1\,\text{м/с}$.
}

\end{document}
